\chapter{O Mercado de TI Segundo o Stack Overflow Developer Survey de 2021}

\section{Sobre a pesquisa}

Para complementar o panorama traçado a partir da pesquisa de 2019, vale a pena revisitar os resultados do Stack Overflow Developer Survey de 2021, realizado em maio daquele ano, com a participação de cerca de 83 mil desenvolvedores de 181 países ao redor do mundo (foram descartadas respostas consideradas não confiáveis, por exemplo aquelas cujo tempo de preenchimento foi curto demais para refletir respostas ponderadas). A pesquisa é dividida em várias seções: perfil do desenvolvedor, tecnologia, trabalho, comunidade e metodologia. Este capítulo percorre os principais achados de cada uma delas, permitindo comparar a evolução do mercado entre 2019 e 2021. De maneira geral, quase 60\% das pessoas que responderam à pesquisa aprendem a programar a partir de recursos online. Um padrão interessante que se repete: desenvolvedores mais jovens tendem a aprender a partir de cursos, fóruns e outros recursos online, enquanto desenvolvedores de idade mais avançada tendem a aprender por meios mais tradicionais, como ensino formal em escolas e livros.

\section{Perfil do desenvolvedor}

O Stack Overflow funciona como uma comunidade verdadeiramente global: a pesquisa recebeu respostas de praticamente todos os países do mundo. Os Estados Unidos e a Índia continuam a fornecer o maior volume de respostas, seguidos por Alemanha e Reino Unido. O Brasil correspondeu a 2,7\% do total de aproximadamente 49 mil respostas geográficas coletadas. Alguns dados demográficos e de formação relevantes:

\begin{itemize}
  \item A grande maioria dos desenvolvedores escreveu sua primeira linha de código entre 11 e 17 anos de idade;
  \item 59,53\% aprenderam a programar a partir de recursos online (blogs, vídeos etc.), 53,59\% em escolas e universidades e 51,53\% a partir de livros ou meios físicos — ou seja, é comum combinar mais de uma fonte de aprendizado;
  \item Considerando anos totais de experiência com programação (incluindo tempo de formação), 29,91\% das pessoas relataram entre 5 e 9 anos de experiência; no Brasil, especificamente, a média foi de cerca de 11,5 anos;
  \item Considerando apenas experiência profissional (excluindo educação), 31,26\% relataram entre 1 e 4 anos de experiência no mercado de trabalho;
\end{itemize}

\begin{itemize}
  \item Quanto à formação: 42,37\% possuem nível universitário (bacharelado), 20,99\% possuem título de mestre e apenas 3\% possuem doutorado. Isso indica que, mesmo em países como o Brasil, onde a profissão não é regulamentada e não exige diploma, o caminho acadêmico continua sendo comum entre os profissionais da área;
  \item Quanto à idade: 39,52\% têm entre 25 e 34 anos e 25,47\% têm entre 18 e 24 anos — ou seja, cerca de 65\% dos desenvolvedores têm entre 18 e 34 anos, confirmando uma profissão de perfil etário relativamente jovem;
  \item Quanto a gênero: 91,67\% se identificaram como homens, 5,31\% como mulheres e cerca de 1\% como pessoas transgênero, evidenciando uma desproporção de gênero significativa na área. Um padrão interessante relacionado à experiência por cargo: executivos e gerentes tendem a ter mais anos de experiência (na casa de 13 a 16 anos, em média), enquanto profissionais de ciência de dados e aprendizado de máquina apresentam, em média, menos anos de experiência do que mesmo pesquisadores acadêmicos. Entre os papéis específicos, DevOps aparece com média de 11,26 anos de experiência, e desenvolvedores mobile e front-end com cerca de 9 anos. Quanto ao tipo de atuação declarada (a pergunta permitia múltiplas respostas), 49\% se identificaram como desenvolvedores full stack, 43\% como back-end, 27\% como front-end, 16\% como desenvolvedores desktop e 14\% como desenvolvedores mobile. Vale notar que o papel de designer caiu significativamente em relação a 2020, sendo ultrapassado pelo de administrador de sistemas.
\end{itemize}

\section{Tecnologias mais usadas, mais amadas e mais temidas}

Um dos aspectos mais interessantes da pesquisa é a comparação entre tecnologias mais usadas, mais adoradas (loved), mais temidas (dreaded) e mais desejadas (wanted) para o futuro. Essa distinção é importante: uma tecnologia pode ser muito usada no mercado sem ser, necessariamente, a preferida dos desenvolvedores.

\begin{itemize}
  \item Linguagens mais usadas em 2021: JavaScript continuou na liderança pelo nono ano consecutivo (usada por 64,96\% dos profissionais respondentes), seguida por HTML/CSS (56\%), Python (48\%) e SQL, que ultrapassou o Java e se tornou a terceira linguagem mais popular entre desenvolvedores web;
  \item Bancos de dados: MySQL (50\%) permanece como o mais usado, seguido por PostgreSQL (40,42\%), Microsoft SQL Server (32,18\%) e MongoDB (27,7\%) — uma mistura de bancos relacionais tradicionais e bancos NoSQL;
  \item Plataformas de nuvem: AWS lidera com 54,22\% de preferência, seguida por Google Cloud Platform (31\%) e Microsoft Azure (30,77\%);
  \item Frameworks web: Node.js aparece com 40,14\%, seguido por React (34,42\%), Express (23,02\%) e Angular (22\%) — o React ultrapassou o jQuery como framework web mais usado;
\end{itemize}

\begin{itemize}
  \item Ferramentas: cerca de 90\% dos desenvolvedores usam Git como sistema de controle de versão, tornando-o praticamente indispensável na profissão; 48,85\% relataram uso extensivo do Docker no último ano;
\end{itemize}

\begin{itemize}
  \item IDEs: o Visual Studio Code lidera disparadamente, usado por 71,06\% dos desenvolvedores, seguido por IntelliJ IDEA (28,74\%), Eclipse (15,87\%) e PyCharm (7,47\%).
\end{itemize}

Quanto às tecnologias mais amadas e mais temidas:

\begin{itemize}
  \item Rust é a linguagem mais adorada (87\% de aprovação entre quem a utiliza), mas também considerada difícil por 13,02\% dos respondentes;
\end{itemize}

\begin{itemize}
  \item TypeScript é amado por 72,73\% e temido por 27,27\%;
\end{itemize}

\begin{itemize}
  \item Java é amado por 47,15\%, mas considerado difícil por 52,82\% — uma das linguagens mais temidas apesar de ser extremamente comum em ambientes universitários;
\end{itemize}

\begin{itemize}
  \item Python é uma das linguagens que as pessoas mais desejam aprender no futuro, ao passo que apenas 6\% desejam aprender Java na sequência da carreira, e 15,29\% desejam aprender TypeScript;
\end{itemize}

\begin{itemize}
  \item Entre bancos de dados, o Redis é o mais amado (70,71\%), enquanto o MongoDB, apesar de popular (51\% de uso), é considerado intimidador por 48,65\% dos respondentes;
\end{itemize}

\begin{itemize}
  \item Entre frameworks web, React tem cerca de 72\% de preferência, e o jQuery, apesar de tradicional, vem perdendo espaço frente a frameworks mais modernos.
\end{itemize}

\section{Salários por tecnologia e por cargo}

A pesquisa também traçou um panorama salarial detalhado. Entre linguagens de programação, o Clojure pagava, em 2021, o maior salário mediano anual (cerca de US\$ 95.000), seguido por linguagens como Erlang, F\#, Scala, Go e Ruby — muitas delas associadas a nichos especializados do mercado, o que explica salários mais altos. PHP, em contraste, aparece entre as linguagens de menor remuneração mediana (cerca de US\$ 38.000 anuais nos dados apresentados), apesar de ser amplamente utilizada. Quanto a cargos nos Estados Unidos, executivos de nível sênior lideram com média de US\$ 177.000 anuais, seguidos por gerentes de engenharia, engenheiros de confiabilidade de sites (SRE) e especialistas em DevOps (cerca de US\$ 135.000). Desenvolvedores back-end, de jogos, full stack e front-end aparecem em faixas entre US\$ 115.000 e US\$ 133.000 anuais nos Estados Unidos. É importante lembrar que esses valores refletem o mercado americano; a mediana global, considerando todos os países da pesquisa, é substancialmente menor. Um padrão consistente aparece quando cruzamos salário com anos de experiência: cargos de gestão têm, em média, tanto salários mais altos quanto mais anos de experiência acumulada, e a experiência profissional tende a se correlacionar positivamente com a remuneração em praticamente todos os papéis analisados.

\section{Emprego, empresas e a comunidade Stack Overflow}

Sobre o tipo de vínculo empregatício, 81\% dos desenvolvedores profissionais trabalham em regime integral (full-time). O percentual de profissionais que trabalham como autônomos (freelancers) cresceu de 9,5\% em 2020 para 11,2\% em 2021, sugerindo uma tendência de crescimento do trabalho independente na área. Quanto ao tamanho das empresas, cerca de 40\% dos desenvolvedores trabalham em empresas de 20 a 500 funcionários, e apenas 6,16\% em empresas de 500 a 1.000 funcionários. Quanto ao uso e à participação na comunidade do Stack Overflow:

\begin{itemize}
  \item Praticamente todos os desenvolvedores visitam o site quando encontram algum problema;
\end{itemize}

\begin{itemize}
  \item 23,76\% visitam o site várias vezes ao dia e cerca de 30\% quase diariamente;
\end{itemize}

\begin{itemize}
  \item 82,37\% possuem conta no Stack Overflow, mas 45,56\% participam (perguntando ou respondendo) menos de uma vez por mês;
\end{itemize}

\begin{itemize}
  \item Apenas 28\% se consideram parte da comunidade do Stack Overflow de forma efetiva.
\end{itemize}

Uma seção curiosa da pesquisa trata de como os desenvolvedores resolvem problemas no dia a dia: quase 90\% recorrem ao Google, 80\% visitam o Stack Overflow, 44\% assistem a tutoriais, e uma fração relevante (37\%) relata sair para caminhar ou praticar alguma atividade física antes de retomar o problema — um lembrete de que descanso e distância do problema muitas vezes ajudam a encontrar soluções mais rapidamente do que a insistência imediata.

\begin{infobox}{Dica}
Vale registrar esse achado da pesquisa como um conselho prático: quando você travar em um problema de código e sentir frustração crescente, uma pausa breve — uma caminhada, uma xícara de café, uma troca de ambiente — é uma estratégia legítima e frequentemente eficaz, não uma forma de procrastinação.
\end{infobox}

Síntese do Capítulo

\begin{itemize}
  \item A pesquisa de 2021 reuniu quase 83 mil respostas de 181 países, confirmando o Stack Overflow como referência global para entender o mercado de TI.
\end{itemize}

\begin{itemize}
  \item JavaScript segue como a linguagem mais usada pelo nono ano consecutivo, mas Python ultrapassou SQL em popularidade geral e Rust é a linguagem mais amada pelos desenvolvedores.
\end{itemize}

\begin{itemize}
  \item O Visual Studio Code e o Git são, respectivamente, o ambiente de desenvolvimento e a ferramenta de controle de versão mais difundidos no mercado.
\end{itemize}

\begin{itemize}
  \item Salários variam fortemente por tecnologia (Clojure e Ruby entre os mais bem pagos, PHP entre os menos) e por país, sendo os valores dos EUA muito superiores à mediana global.
\end{itemize}

\begin{itemize}
  \item O trabalho autônomo (freelance) cresceu de 9,5\% para 11,2\% entre 2020 e 2021, indicando uma tendência de maior flexibilidade de vínculo na área.
\end{itemize}

\begin{itemize}
  \item A maioria dos desenvolvedores tem entre 18 e 34 anos e formação superior, ainda que a profissão nem sempre exija diploma formal.
\end{itemize}

\begin{itemize}
  \item Resolver problemas de programação envolve tanto pesquisa ativa (Google, Stack Overflow, tutoriais) quanto pausas estratégicas, segundo o próprio relato dos desenvolvedores.
\end{itemize}
